% Plantilla creada por Francisco Javier Vázquez Tavares
% Esta platilla se creo para facilitar la creación de documentos para tareas.

\documentclass{tareaClass}
%\selectlanguage{spanish} 
%\usepackage[spanish,onelanguage]{algorithm2e} %for psuedo code
\usepackage{listings}
%\usepackage{enumitem}

% For newtcbtheorem
\usepackage[most]{tcolorbox}
\usepackage{cleveref}
\usepackage{fancybox}		% Recuadros emph eqn

\usepackage{empheq}			% Recuadros para ecuaciones


\newtcbtheorem[]{prob}{Problem}%
    {enhanced,
    colback = black!5, %white,
    colbacktitle = black!5,
    coltitle = black,
    boxrule = 0pt,
    frame hidden,
    borderline west = {0.5mm}{0.0mm}{black},
    fonttitle = \bfseries\sffamily,
    breakable,
    before skip = 3ex,
    after skip = 3ex
}{def}

\materia{Algebraic and Transcendental Functions}
\profe{Professor: Fco. Vazquez}
\title{ Procedures }

\integrantes{~}
\matriculas{~}

\begin{document}

\lhead{Algebraic and Transcendental Functions}
\rhead{\today}

\maketitle

\section{Introduction}

This document explains how to do some of the exercises that you have done during the course.
The idea is to help you get better at maths by practising the steps.
I hope this helps you.
This document is going to keep changing, so please be patient.

\subfile{docs/hwk1.tex}
%\subfile{docs/hwk2.tex}
%\subfile{docs/hwk3.tex}
%\subfile{docs/hwk4.tex}
%\subfile{docs/hwk5.tex}
%\subfile{docs/hwk6.tex}


%\subfile{docs/quiz.tex}


\end{document}
